\documentclass{report}
\usepackage[utf8]{inputenc}
\usepackage[T1]{fontenc}
\usepackage{geometry}
\usepackage{graphicx}
\usepackage{float}
\usepackage{hyperref}
\usepackage[english,indonesian]{babel}

\title{Optimasi Pemrosesan Template LaTeX Menggunakan Pendekatan Berbasis AI}
\author{Academic Support Agent}
\date{\today}

\begin{document}
\maketitle

\begin{abstract}

\end{abstract}

\chapter{PENDAHULUAN}


\chapter{TINJAUAN PUSTAKA}


% Original content (commented out):
% %Modelo desenvolvido pelo discente Bruno Sávio Fichos do curso de Engenharia Mecânica do CEFET-MG
%Versao 2021.1

\documentclass[12pt,oneside, bibliography=totoc]{report}
\usepackage{definitions}

\pagestyle{plain}
\pagenumbering{gobble}
%%%%%%%%%%%%%%%%%%%%%%%%%%%%%%%%%%%%%%%%%%%%%%%%%%%%%%%%%%%%%
\makeglossaries
\input{listas/abreviacoes.tex}

\makenomenclature

\begin{document}
    
    \input{capa.tex}
    \input{subcapa.tex}
    %\input{termo_aprovacao.tex}
    %\include{tcc1/parte1}
    %\include{tcc1/parte2}
    %\include{tcc1/parte3}
    %\input{dedicatoria.tex}
    %\input{agradecimentos.tex}
    
    \clearpage
    \input{capitulos/resumo.tex}
    \input{capitulos/abstract.tex}
    \listoffigures
    \pagebreak
    \listoftables
    \printglossary[type=\acronymtype]
    \printglossary
    \printnomenclature
    \clearpage
    
    \tableofcontents{}
    \clearpage
    \setcounter{page}{1}
    \pagenumbering{arabic}
    \include{capitulos/auxiliar}
    \include{capitulos/introducao}
    \include{capitulos/justificativa}
    \include{capitulos/objetivos}
    \include{capitulos/revisao_bibliografica}
    \include{capitulos/metodologia}
    \include{capitulos/resultados}
    \include{capitulos/conclusão}
    \include{capitulos/propostas_futuros}
    \include{listas/simbolos}
    \include{listas/referencias}
    
    
    
\end{document}


\end{document}